% DFS 判环四种情形：(a) 有向 1->2->3，检查 3->1 时 1 在栈内；(b) 有向 1->2, 1->3, 3->2，检查 3->2 时 2 已完成；
% (c) 无向路径 1-2-3，回看父边不算环；(d) 无向三角形，3 看到已发现且非父点 1。
\begin{tikzpicture}[font=\figurefont\fontsize{8}{10}\selectfont,>=Stealth,
  vtx/.style={circle,draw=FigureStroke,line width=.8pt,minimum size=6.2mm,inner sep=0pt,text=FigureInk,fill=white},
  stk/.style={vtx,fill=FigureYellowFill},
  fin/.style={vtx,fill=FigureBlueFill},
  tree/.style={draw=FigureBlue,line width=1.2pt},
  back/.style={draw=FigureRed,line width=1.2pt},
  other/.style={draw=FigureMuted,line width=.8pt,dashed},
  note/.style={font=\FigureNoteFont\fontsize{7.5}{9}\selectfont,text=FigureMuted,align=center}]
% (a)
\begin{scope}
  \node[stk] (a1) at (0,1.1) {1}; \node[stk] (a2) at (-0.6,0) {2}; \node[stk] (a3) at (0.6,0) {3};
  \draw[tree,->] (a1)--(a2); \draw[tree,->] (a2)--(a3);
  \draw[back,->] (a3)--(a1);
  \node[note] at (0,-0.75) {$3\to1$指向栈内点\\沿父指针恢复环};
  \node at (0,-1.6) {(a) 有向回边};
\end{scope}
% (b)
\begin{scope}[xshift=2.65cm]
  \node[stk] (b1) at (0,1.1) {1}; \node[fin] (b2) at (-0.6,0) {2}; \node[stk] (b3) at (0.6,0) {3};
  \draw[tree,->] (b1)--(b2); \draw[tree,->] (b1)--(b3);
  \draw[other,->] (b3)--(b2);
  \node[note] at (0,-0.75) {$3\to2$指向已完成点\\不构成环};
  \node at (0,-1.6) {(b) 指向已完成点};
\end{scope}
% (c)
\begin{scope}[xshift=5.3cm]
  \node[stk] (c1) at (0,1.1) {1}; \node[stk] (c2) at (-0.6,0) {2}; \node[vtx] (c3) at (0.6,0) {3};
  \draw[tree] (c1)--(c2);
  \draw[FigureStroke,line width=.8pt] (c2)--(c3);
  \draw[other,->,bend left=35] (c2) to (c1);
  \node[note] at (0,-0.75) {从2回看父边$\{1,2\}$\\跳过，不算环};
  \node at (0,-1.6) {(c) 无向父边};
\end{scope}
% (d)
\begin{scope}[xshift=7.95cm]
  \node[stk] (d1) at (0,1.1) {1}; \node[stk] (d2) at (-0.6,0) {2}; \node[stk] (d3) at (0.6,0) {3};
  \draw[tree] (d1)--(d2); \draw[tree] (d2)--(d3);
  \draw[back] (d3)--(d1);
  \node[note] at (0,-0.75) {3遇到已发现的\\非父点1，成环};
  \node at (0,-1.6) {(d) 无向非父边};
\end{scope}
\end{tikzpicture}
